\DocumentMetadata{
  pdfversion=2.0,
  pdfstandard=ua-2,
  lang=en-US,
  tagging=on,
  tagging-setup={math/setup=mathml-SE,tabsorder=structure}
}
\documentclass[11pt,letterpaper]{article}
\usepackage[margin=0.68in,includefoot]{geometry}
\usepackage{newcomputermodern}
\usepackage{amsmath,amscd,mathtools}
\providecommand{\square}{\mdlgwhtsquare}
\usepackage[hidelinks]{hyperref}
\hypersetup{
  pdftitle={Partial Differential Equations PhD Exam, January 2014},
  pdfauthor={University of Florida Department of Mathematics},
  pdfsubject={Graduate mathematics examination},
  pdfdisplaydoctitle=true,
  bookmarksopen=true
}
\setcounter{secnumdepth}{0}
\setlength{\parindent}{0pt}
\setlength{\parskip}{0.28em}
\setlength{\emergencystretch}{3em}
\setlength{\leftmargini}{2.15em}
\setlength{\leftmarginii}{2.65em}
\setlength{\leftmarginiii}{3em}
\renewcommand{\labelenumi}{\textbf{\theenumi.}}
\pagestyle{empty}
\raggedbottom
\allowdisplaybreaks
\newenvironment{examproblems}[1]{%
  \begin{enumerate}%
  \setcounter{enumi}{#1}%
  \setlength{\topsep}{0.35em}%
  \setlength{\partopsep}{0pt}%
  \setlength{\itemsep}{0.48em}%
  \setlength{\parsep}{0.18em}%
}{\end{enumerate}}
\newenvironment{examsubparts}{%
  \begin{enumerate}%
  \setlength{\topsep}{0.18em}%
  \setlength{\partopsep}{0pt}%
  \setlength{\itemsep}{0.16em}%
  \setlength{\parsep}{0.1em}%
}{\end{enumerate}}
\newenvironment{examinstructions}{%
  \par\begingroup\itshape
}{%
  \par\endgroup\vspace{0.18em}
}
\makeatletter
\renewcommand\section{\@startsection{section}{1}{0pt}%
  {-1.25ex plus -.25ex minus -.1ex}%
  {0.55ex plus .1ex}%
  {\normalfont\large\bfseries}}
\renewcommand{\@maketitle}{%
  \newpage\null\vskip -1.2em%
  {\footnotesize\bfseries
    \href{https://www.ufl.edu/}{University of Florida}, \href{https://math.ufl.edu/}{Department of Mathematics}\par}%
  \vskip 0.18em%
  \tagstructbegin{tag=Title}%
    {\LARGE\bfseries\href{https://gma.math.ufl.edu/exams/pde-phd/}{Partial Differential Equations PhD Exam}, January 2014\par}%
  \tagstructend%
  \vskip 0.35em%
  \tagmcbegin{artifact}%
    \hrule height 1pt%
  \tagmcend%
  \vskip 0.55em%
}
\makeatother
\title{Partial Differential Equations PhD Exam, January 2014}
\author{}
\date{}
\begin{document}
\maketitle
\thispagestyle{empty}
\begin{examproblems}{0}
\item
Let~\(\Omega\) be an open subset of \(\mathbb{R}^n\) and let \(u\in C^2(\Omega)\) be a harmonic function in~\(\Omega\). Prove that~\(u\) is analytic in~\(\Omega\).
\item
Let~\(U\) be an open and bounded set in \(\mathbb{R}^n\) with smooth boundary. Suppose \(u_1\) and \(u_2\) are both smooth solutions to 
\[
\begin{cases}
u_t-\Delta u=0, & \text{in } U_T:=U\times(0,T],\\
u=g, & \text{on } \partial U\times[0,T].
\end{cases}
\]
 If \(u_1(x,T)=u_2(x,T)\) for all \(x\in U\), prove that \(u_1\equiv u_2\) within \(U_T\).
\item
Suppose~\(u\) is a smooth solution to 
\[
u_{tt}-\Delta u=0, \qquad \text{in } \mathbb{R}^n\times(0,\infty)
\]
 and 
\[
u\equiv u_t\equiv 0 \qquad \text{on } B(x_0,t_0)\times\{t=0\}.
\]
 Let 
\[
C=\{(x,t)\mid 0\leq t\leq t_0,\ |x-x_0|\leq t_0-t\}.
\]
 Prove that \(u\equiv 0\) within~\(C\).
\item
Use characteristics to solve the following equation: 
\[
\begin{cases}
x_1u_{x_1}+2x_2u_{x_2}+u_{x_3}=3u, & \mathbb{R}_+^3,\\
u(x_1,x_2,0)=g(x_1,x_2), & \partial\mathbb{R}_+^3,
\end{cases}
\]
 where~\(g\) is a given smooth function.
\item
Fix \(\alpha>0\) and let \(U=B(0,1)\) (open ball). Show there exists a constant~\(C\), depending only on~\(n\) and~\(\alpha\), such that 
\[
\int_U u^2\,dx\leq C\int_U |Du|^2\,dx,
\]
 provided 
\[
|\{x\in U\mid u(x)=0\}|\geq\alpha, \qquad u\in H^1(U).
\]
\item
Use Hopf–Lax formula to find a solution of 
\[
\begin{cases}
u_t+\dfrac{1}{2}|Du|^2=0, & \mathbb{R}^n\times(0,\infty),\\
u=-\dfrac{1}{2}|x|, & \mathbb{R}^n\times\{t=0\}.
\end{cases}
\]
 Is this solution unique in the weak sense?
\item
Let~\(U\) be a bounded subset of \(\mathbb{R}^n\), \(u\in W^{2,p}(U)\cap W_0^{1,p}(U)\), and \(2\leq p<\infty\). Prove that 
\[
\int_U |Du|^p\,dx\leq C\left(\int_U |u|^p\,dx\right)^{\frac12}\left(\int_U |D^2u|^p\,dx\right)^{\frac12}
\]
 for some~\(C\) independent of~\(u\).
\end{examproblems}
\end{document}
